\chapter{Discussion, Conclusion, and Future Work}
\label{chap:discussion}

\section{Discussion of Research Questions}

\subsection{RQ1: Complexity Classification and Routing}

The implementation demonstrates that a business-workflow QA system can expose
query-complexity classification as a runtime deployment rather than a hidden
hard-coded heuristic. The classifier result carries probabilities, confidence,
margin, latency, configured deployment, actual deployment, and fallback status.
Adaptive routing consumes the scored result, while explicit modes permit
controlled route experiments.

RQ1 is answered for the prepared source-grouped validation split. DistilBERT
and T5-small both achieved \FinalDistilBERTAccuracy{} accuracy and
\FinalDistilBERTMacroFOne{} macro F1, with perfect advanced-class
precision/recall/F1. Their only errors were six complex questions routed as
moderate, producing \FinalUnderRouting{} under-routing and no over-routing.
DistilBERT was substantially faster at \FinalDistilBERTLatency{} versus
\FinalTFiveLatency{}, while T5-small was better calibrated
(\FinalTFiveECE{} versus \FinalDistilBERTECE{} expected calibration error).
The models are therefore equal on held-out predictive quality but differ in
the latency--calibration trade-off.

Adding \texttt{advanced} creates a clearer ownership boundary. L3 remains
decomposition plus multi-query evidence aggregation. L4 owns bounded iterative
planning and tools. This distinction prevents every difficult question from
being called ``agentic.'' It also creates a more demanding classification
problem because complex and advanced questions can share surface features.
Outcome-derived labels and route-cost regret are therefore more appropriate
than a purely linguistic labeling rule.

\subsection{RQ2: Groundedness, Quality, and Efficiency}

The system can execute and compare direct, fixed-retrieval, decomposed, and
agentic configurations using the same orchestration service. Query embeddings
are tied to active indexes, citations are validated, remote data egress is
explicit, and every model attempt records tokens, cost, and latency. These are
necessary conditions for a meaningful accuracy--efficiency comparison.

The completed configuration experiment provides primary current evidence for
RQ2. Adaptive with DistilBERT ranked first at quality
\DistilConfigQuality{}, followed by Fixed L2 at \LTwoConfigQuality{}, Adaptive
with T5-small at \TFiveConfigQuality{}, and Fixed L3 at
\LThreeConfigQuality{}. The DistilBERT configuration averaged
\DistilConfigLatency{} per case; the T5-small configuration averaged
\TFiveConfigLatency{}, demonstrating that router latency can materially affect
end-to-end responsiveness even when both classifiers predict the same route.

This result is evidence that Adaptive routing is competitive in the tested
configuration matrix, not proof of general superiority. The experiment used
only \PrimaryExperimentCases{} compatible cases from a
\PrimaryExperimentKBDocuments-document KB, and both Adaptive classifiers routed
every case to L2. Furthermore, top-\textit{k} and generator settings differ
between configurations. A causal route comparison therefore requires a larger
compatible sample and one-factor-at-a-time controls.

\subsection{RQ3: Explainable Evaluation}

RAGXplain is operationally integrated rather than represented by a mock report.
It receives real questions, contexts, answers, ground truths, configuration
snapshots, and judge deployment. Its executive summary and protocols load
automatically in the application. Diagnosis completed for all twelve runs and
consistently identified retrieval breadth, evidence coverage, and generator
grounding as the primary constraints. Recommended changes correspond to actual
fields: \textit{top-k}, reranker, structure-aware chunking, query expansion,
grounding prompt, citations, and response structure.

This supports the mechanism proposed by RQ3: a metric can be translated into a
specific configuration hypothesis. It does not yet establish that applying the
recommendation causes a statistically reliable improvement. The required next
experiment is a closed loop: preserve the original configuration, clone and
change only the recommended factors, rerun the same case sample, and compare
quality, latency, and cost. RAGXplain should be considered decision support,
not an autonomous optimizer.

\section{Engineering Trade-offs}

\subsection{Quality, Latency, and Cost}

Increasing \textit{top-k} may improve recall but consumes context tokens and can
introduce distractors. Reranking can recover precision but adds another model
call. L3 can retrieve missing evidence but multiplies searches. L4 can recover
from uncertain plans but has the largest latency and failure surface. The
Adaptive approach is valuable only when routing savings and quality gains
outweigh classifier overhead and misrouting cost.

The AI Model Gateway makes those trade-offs measurable. Local deterministic
models support inexpensive functional tests. OpenRouter supports
experimentation. Direct OpenAI or Gemini connections provide stable provider
paths. Ollama and vLLM support self-hosted deployments. Fallback improves
availability, but it is intentionally blocked after streamed output begins so
one answer never combines providers invisibly.

\subsection{Observability and Privacy}

Full payload traces accelerate diagnosis because they preserve prompts,
retrieval candidates, scores, model parameters, responses, and usage. The same
payloads can contain sensitive organization content. The prototype mitigates
this through JWT access, redaction, retention, compression, and size limits,
but a production deployment would require tenant isolation, finer-grained trace
authorization, encrypted artifact storage, audit logs, retention policy
approval, and data-loss prevention.

\subsection{Flexibility and Complexity}

The application exposes a large configuration surface because the capstone
studies RAG component choices. This is useful for researchers but can overwhelm
ordinary users. Saved configurations, provider capability filtering, inherited
embeddings, presets, and collapsible sections reduce that complexity. A
production product should separate an administrator/research workspace from a
minimal end-user chat.

\section{Limitations and Threats to Validity}

\begin{itemize}
  \item \textbf{Small compatible sample.} The primary experiment contains
  \PrimaryExperimentCases{} cases: \PrimarySimulatedCases{} Simulated,
  \PrimarySyntheticCases{} Synthetic, and \PrimaryExpertCases{}
  ExpertWritten. This is insufficient for significance or confidence-interval
  claims.
  \item \textbf{Partial knowledge corpus.} The experiment indexed
  \PrimaryExperimentKBDocuments{} of \PreparedCorpusDocuments{} prepared WixQA
  documents. Results apply only to questions compatible with those records.
  \item \textbf{Configuration confounding.} Top-\textit{k} differs between
  configurations, and Fixed L3 uses a local extractive generator while the
  other configurations use GPT-4.1-mini. The leaderboard does not isolate a
  single causal component.
  \item \textbf{Limited route coverage.} Both Adaptive classifiers routed all
  twenty compatible cases to L2. The experiment therefore does not validate
  Adaptive L1, L3, or L4 selection behavior.
  \item \textbf{Judge dependence.} Factuality, context recall, and RAGXplain
  metrics depend on an LLM judge and may vary by model, prompt, and provider.
  \item \textbf{Template-label bias.} The balanced four-class dataset includes
  \PreparedTemplateRecords{} template examples. Although source groups are
  disjoint, templates may be easier to classify than authentic questions.
  \item \textbf{Domain limitation.} Wix support content is enterprise-like but
  represents one product ecosystem and language distribution.
  \item \textbf{Prototype operations.} The system is tested primarily on a
  Windows workstation and local PostgreSQL. It is not a high-availability,
  multi-tenant production deployment.
  \item \textbf{L4 maturity.} Tool availability is intentionally limited.
  Unavailable Drive, OneDrive, and database connectors constrain evaluation of
  broad research tasks.
\end{itemize}

\section{Future Work}

\subsection{Immediate Experimental Work}

The next priority is to extend the completed experiment into a controlled
matrix over a larger compatible sample and, when resources permit, the full
\PreparedCorpusDocuments-document WixQA corpus. Each comparison should change
one factor while preserving case IDs, KB index version, generator, prompts,
judge, and unrelated settings. Required experiments include:

\begin{itemize}
  \item L4 Advanced RAG against matched Adaptive, L2, and L3 controls;
  \item BM25, dense, and hybrid retrieval with the same generator and cases;
  \item reranking disabled and enabled over the same first-stage candidates;
  \item top-\textit{k} values 2, 5, and 7; and
  \item baseline versus RAGXplain-recommended prompt and retrieval settings.
\end{itemize}

Classifier confidence and margin thresholds should be selected on validation
data and frozen before the expanded route benchmark. Repeated or bootstrapped
runs should report uncertainty in addition to means.

\subsection{Model and Retrieval Enhancements}

Future work can add cross-encoder rerankers, multilingual embeddings,
query-expansion models, learned hybrid weights, and answer-confidence models.
Index-level experiments should compare fixed, structure-aware, semantic, and
parent-child chunking while holding generator and case sample constant.

The classifier could use outcome-based labels generated from repeated route
trials rather than one deterministic snapshot. Cost-aware training could
penalize under-routing more strongly than over-routing while still discouraging
unnecessary L4 execution.

\subsection{Connectors and L4 Tools}

Google Drive, OneDrive, SharePoint, and database ingestion require production
OAuth, incremental synchronization, deletion propagation, permission mapping,
and connector-specific rate limits. Only after those controls exist should the
tools become visible to L4 planners. Public web retrieval should gain domain
allowlists, content provenance, and stronger prompt-injection filtering.

\subsection{Production Hardening}

Production work includes external identity integration, multi-tenant
authorization, managed secret storage, artifact encryption, distributed
cancellation, horizontally scalable workers, object storage, backups,
disaster recovery, rate limits, formal threat modeling, and load tests.
Telemetry can later be exported to OpenTelemetry while preserving current
trace/span identifiers.

\subsection{Human Evaluation}

Automated metrics should be complemented by workflow experts who assess
correctness, completeness, actionability, source support, and harmful
omissions. The implemented answer-version feedback model provides a starting
point, but a controlled human study should randomize configurations, blind
reviewers where possible, and measure agreement.

\section{Conclusion}

This capstone transformed the proposal for adaptive business-workflow QA into a
working full-stack research system. \SystemName{} connects knowledge
processing, versioned embedding indexes, hybrid retrieval, four-level adaptive
routing, local and hosted models, streaming interaction, citations, durable
traces, configuration-based evaluation, RAGXplain diagnostics, and analytics.
Its architecture is designed so each component can be independently tested and
its runtime behavior can be reconstructed.

The implementation evidence shows that the complete experimental workflow is
operational, both four-class classifiers achieve strong grouped-validation
performance, and Adaptive with DistilBERT is the highest-ranked configuration
in the current 20-case experiment. It also shows why disciplined evaluation
matters: corpus compatibility and controlled configuration snapshots are
preconditions for valid route comparison, and explainable metrics are most
useful when they map to concrete configuration changes.

The project does not claim statistically significant Adaptive RAG superiority
from the current sample. Its contribution is a transparent, reproducible
system with trained routers, persisted configuration experiments, explainable
diagnosis, and sufficient observability to run and defend the larger controlled
matrix without rebuilding the system.
